\begin{document}
\frontmatter
\begin{titlepage}
\centering
\vspace*{26mm}
{\Large THE MATHEMATICAL FOUNDATIONS\par}
\vspace{4mm}
{\Large OF THE UNNS SUBSTRATE\par}
\vspace{15mm}
{\Huge\bfseries Volume II\par}
\vspace{5mm}
{\LARGE Collective Dynamics, Structural Rigidity,\par and Effective Predictability\par}
\vspace{10mm}
{\large From Generative Multiplicity to Exact Scaling\par and Arithmetic Certification\par}
\vfill
{\large Working manuscript\par}
{\large 23 October 2026\par}
\vspace{6mm}
{\normalsize Source baseline: GitHub reference collection\par commit \texttt{bb8b12d753c600c597feb00df17fe139c7cc65fa}\par}
\end{titlepage}

\chapter*{Abstract}
\addcontentsline{toc}{chapter}{Abstract}
Volume I determined what the UNNS Substrate generates: which nested occurrences realise a given value, how their routes meet again, how rigidly they are distinguished, and why no interaction law is forced by the arithmetic alone. This volume asks what happens when generative structures evolve \emph{collectively}. We fix one explicit interaction rule, the collective map \(C(x)=x+L(x)^2/G(x)\) on the positive rationals, where \(L\) and \(G\) are the least common multiple and the greatest common divisor of the moduli carried by the occurrence fibre of \(x\). Once the rule is selected, its dynamics are governed by theorems, and five of them organise the volume.

\emph{Unbounded generative multiplicity.} Along every orbit the fibre size \(R(x_n)\) is non-decreasing and tends to infinity, the gcd \(G(x_n)\) can only decrease through divisors, and the reach set, although infinite, has logarithmic counting function and density zero. \emph{Eventual exact scaling.} Every orbit satisfies \(x_{n+P}=\rho^2x_n\) from some time on. The proof passes through an arithmetic state description in which a persistent orbit is a finite-state process driven by a single integer coordinate \(\gamma\) and residue classes of prime exponents; the persistent regimes with \(G_\infty>1\) exist for every \(g\ge5\) prime to 6 and for several \(g\) divisible by 3, \(G_\infty\) never divides 24, and the eight-step orbit at \(g=11\) shows that no universal three-step law holds. \emph{Universal single-prime rigidity.} The carries of a persistent orbit form a system of independent cyclic counters read by a scheduler that always fires the smallest pending carry; exactly one prime carries infinitely often, and the scaling factor is a prime power \(\rho=p^s\). \emph{Arithmetic survivor selection.} In a system of independent odometers with fixed tables, the survivor is the prime whose worst-case carry is smallest; a quotient-residue formula computes that carry from finite residue data; and, at a reset state of an orbit with prime \(g\) that satisfies the stated certificate hypotheses, a finite residue test decides the survivor and fixes the entire future. For an arbitrary initial state the survivor is determined by running the exact carry-level predictor to a reset state, within the transient bounds below. \emph{Effective finite predictability.} From the initial pair \((g_0,\gamma_0)\) alone, an explicit and computable bound limits the number of carry events before the orbit enters its final cycle; the reset-descent part of this bound, \(\kappa(\gamma_0)\), is attained at each level \(\kappa=1,\dots,11\) by explicit witnesses.

The volume develops all seven stages with complete proofs, states every hypothesis and scope, and supports each statement with exact integer computation reproducible from the accompanying package.

\tableofcontents
\mainmatter

\chapter*{Introduction: from generative structure to predictive dynamics}
\addcontentsline{toc}{chapter}{Introduction: from generative structure to predictive dynamics}

\textbf{Where Volume I ends.} Volume I established the arithmetic of the Substrate. For a positive rational \(x=a/b\) in lowest terms, the occurrences carrying the value \(x\) are indexed by the \emph{moduli}
\[
\Dx(x)=\{d\ge2:\ d^2\mid a,\ d\equiv1\pmod b\},\qquad R(x)=|\Dx(x)|,
\]
(Volume~I, Theorem~1.2). Any two routes through \(x\) meet again exactly on the progression \(x+j\Dcoll\), and for the whole fibre
\[
\Dcoll(x)=\frac{L(x)^2}{G(x)},\qquad L=\lcmop\Dx(x),\quad G=\gcd\{d-1:d\in\Dx(x)\},
\]
with the integral coherence invariant \(\gamma(x)=x/\Dcoll(x)\in\Z_{>0}\) (Volume~I, Theorem~2.5). Volume I also proved that this arithmetic does not select a continuation rule: label-blind admissibility, stationarity and weak label-blindness cannot all hold, and any interaction must be specified. It ended by naming the rule studied here.

\textbf{The collective map.} The map
\[
C(x)=x+\frac{L(x)^2}{G(x)}=x+\Dcoll(x)
\]
sends each value to the first value that \emph{all} routes through it share in common, other than \(x\) itself. It is an explicitly selected interaction rule. Once it is selected, however, what its orbits do is not a further assumption: it is a matter of theorem. Whether the fibre grows, whether the gcd collapses, whether the orbit settles into a rhythm, which primes drive that rhythm, and how long the settling takes are questions with exact answers, and this volume proves them.

\textbf{Five achievements.} The proofs fall into five groups. Each states a different kind of knowledge: that something happens (existence), what constrains it (structure), and how soon (effectiveness).

\begin{center}\small
\renewcommand{\arraystretch}{1.1}
\begin{tabular}{|p{3.9cm}|p{6.7cm}|p{3.3cm}|}
\hline \textbf{Achievement} & \textbf{Statement} & \textbf{Where}\\\hline
Unbounded generative multiplicity & \(R(x_n)\to\infty\) along every orbit; \(D(x_n)\subseteq D(x_{n+1})\); the reach set is sparse & Ch.~1, Thm~\ref{th:F.T19}, \ref{th:F.T31}; Cor.~\ref{cor:logreach}\\\hline
Eventual exact scaling & \(x_{n+P}=\rho^2x_n\) eventually, for every orbit; explicit persistent families & Ch.~3--4, Thm~\ref{th:F.T51}, \ref{th:F.T50}, \ref{th:F.T57}\\\hline
Universal single-prime rigidity & exactly one prime carries infinitely often; \(\rho=p^s\) & Ch.~5, Thm~\ref{th:F.T60}, \ref{th:F.T61}\\\hline
Arithmetic survivor selection & survivor \(=\arg\min\) of the worst-case carry; static certificate for prime \(g\) & Ch.~6, Thm~\ref{th:F.T64}--\ref{th:F.T67}\\\hline
Effective finite predictability & carry events before the final cycle \(\le\Rbd(g_0,\gamma_0)\); reset descent \(\kappa(\gamma_0)\), attained at \(\kappa\le11\) & Ch.~7, Thm~\ref{th:F.T69}--\ref{th:F.T72}\\\hline
\end{tabular}
\end{center}

Existence is the content of Chapters 1--4: eventual arithmetic behaviour is established for every orbit, with explicit regimes showing which behaviours occur. Structure is Chapters 5--6: the carries that drive the behaviour form a system of independent counters under a one-line scheduler, and that system admits exactly one survivor. Effectiveness is Chapter 7: finite bounds, computable from the first arithmetic data \((g_0,\gamma_0)\), control the transient. The volume therefore goes beyond describing patterns or observing regularity numerically; each stage is a theorem with stated hypotheses.

\textbf{Two precise scopes.} Two qualifications govern how the predictive theorems are to be read, and they are stated here once and repeated in the statements of the theorems concerned.

\emph{(i) Carries versus steps.} The principal bound of Theorem~\ref{th:F.T72} counts \emph{carry events}, the steps at which \(L\) changes, not iterations of \(C\). Between carries the orbit runs through plateaus of at most \(4g\) steps (Theorem~\ref{th:F.T37}), so the number of iterations is at most \((\text{carries}+1)\cdot4g+1\) within each segment of constant \(G\), together with the \(G=1\) tail estimate; the theorem records these additional estimates.

\emph{(ii) Prime \(g\) for certification.} The static certificate (Theorem~\ref{th:F.T66}) and its consequence for the eventual regime (Theorem~\ref{th:F.T68}(a)) are proved for prime \(g\), or more generally whenever no prime factor of \(g\) is at most \(m^2\). For composite and even \(g\) the survivor exists and the final cycle is entered within the class-counting bound, but static certification as formulated here is not available at the reset states of the six exact composite and even cycles examined (Theorem~\ref{th:F.T73}); this is open problem O12.2.

\textbf{Reading the volume.} Chapter~\ref{ch:growth} defines the dynamics and proves that the fibre grows without bound. Chapter~\ref{ch:regimes} studies the invariant \(G\), the regime \(G=1\) in closed form, and the transition mechanism between fibres. Chapter~\ref{ch:persist} shows that \(G\) need not collapse, constructs the persistent regimes, and introduces the coordinate \(\gamma\). Chapter~\ref{ch:gamma} proves the one-step law, the eight-step theorem, exact scaling for every orbit, and the carry structure. Chapter~\ref{ch:schedule} derives the scheduler and single-prime dominance. Chapter~\ref{ch:select} develops minimax selection and certification. Chapter~\ref{ch:bounds} proves the effective bounds. The appendices give notation, the complete theorem crosswalk and audit, the algorithms and checks, the mathematical contributions, and the open problems. Every theorem carries a source-record tag (for example F.T51) referring to the FIELD-I records and the 128-record registry; Figure~\ref{fig:arch} shows how the two volumes fit together.

\begin{figure}[t]
\centering
\resizebox{0.98\linewidth}{!}{\begin{tikzpicture}[>=Stealth,font=\scriptsize]
\tikzset{st/.style={draw,rounded corners=2pt,align=center,inner sep=3pt,text width=1.9cm,minimum height=1.0cm}}
\node[st,fill=black!6] (a) at (0,0) {Generation\\ \textit{Vol.~I, 1}};
\node[st,fill=black!6] (b) at (2.5,0) {Reachability\\ \textit{Vol.~I, 4--6}};
\node[st,fill=black!6] (c) at (5.0,0) {Interaction\\ \textit{Vol.~I, 7; II, 1}};
\node[st] (d) at (7.5,0) {Evolution\\ \textit{II, 2--4}};
\node[st] (e) at (10.0,0) {Rigidity\\ \textit{II, 5}};
\node[st] (f) at (12.5,0) {Prediction\\ \textit{II, 6--7}};
\foreach \p/\q in {a/b,b/c,c/d,d/e,e/f} \draw[->] (\p) -- (\q);
\node[lab,below=3pt of a,text width=2.3cm,align=center] {fibres, rendezvous lattices};
\node[lab,below=3pt of b,text width=2.3cm,align=center] {graphs, histories, no-shortcut order};
\node[lab,below=3pt of c,text width=2.3cm,align=center] {nonselection; the map \(C\)};
\node[lab,below=3pt of d,text width=2.3cm,align=center] {monotone orbits, \(\gamma\)-state, exact scaling};
\node[lab,below=3pt of e,text width=2.3cm,align=center] {carry scheduler, one survivor};
\node[lab,below=3pt of f,text width=2.3cm,align=center] {minimax, certificate, effective bounds};
\end{tikzpicture}}
\caption{The two-volume progression. Shaded stages are the content of Volume I; Chapter~1 of this volume begins where the nonselection theorem of Volume I ends.}\label{fig:arch}
\end{figure}

\textbf{Conventions.} All numbers are exact. Rationals are in lowest terms. \(\vv_p\) is the \(p\)-adic valuation, \(\ord_g(p)\) the multiplicative order of \(p\) modulo \(g\), and \(\lcmop\) and \(\gcd\) are as in Volume~I. The orbit of \(x_0\) is \(x_{n+1}=C(x_n)\), and all quantities \(D_n,L_n,G_n,\gamma_n\) are those of \(x_n\). A \emph{scope tag} marks what a statement depends on: scope \textsf{A0} depends on the nest rule alone, scope \textsf{C} on the adoption of \(C\). Every result of this volume other than the two A0 records (F.T17p and F.T30) is scope \textsf{C}. Appendix~\ref{app:notation} fixes the notation and its correspondence with Volume~I and the FIELD-I symbols.
